\chapter{Instalación, ejecución y reproducibilidad}
\label{app:instalacion}

Este anexo documenta cómo reconstruir y ejecutar el sistema de extremo a extremo, en apoyo directo del objetivo de reproducibilidad (OE-08). Todos los comandos se lanzan desde la raíz del repositorio.

\section{Requisitos}
\label{sec:app-requisitos}

El entorno de Python se gestiona con \texttt{uv}, que debe estar instalado en el sistema. Las cargas ligeras y la suite de pruebas corren en cualquier equipo; las cargas pesadas —descarga del corpus, codificación de \emph{embeddings} y llamadas reales al modelo generativo— se ejecutan en el servidor, con la GPU como opción puntual para acelerar la indexación. La generación y el modelo juez requieren \texttt{Ollama} con los modelos locales descargados (generador \texttt{qwen2.5:7b-instruct} y juez \texttt{gemma3:12b}). Para descargar el modelo de \emph{embeddings} se necesita un token de acceso a Hugging Face en la variable de entorno \texttt{HF\_TOKEN}, que nunca se versiona ni se imprime.

\section{Instalación del entorno}
\label{sec:app-instalacion}

\begin{lstlisting}
uv sync    # crea .venv e instala las dependencias fijadas en el lockfile
\end{lstlisting}

No hacen falta ficheros \texttt{.env} ni secretos para las pruebas: la configuración arranca con valores por defecto seguros.

\section{Puerta de calidad (sin red)}
\label{sec:app-calidad}

La suite completa se ejecuta \emph{offline}, sustituyendo modelos y servicios por dobles de prueba:

\begin{lstlisting}
uv run pytest                                          # suite completa, 100% offline
uv run ruff check .                                    # lint (incluye los notebooks)
uv run ruff format .                                   # formato
uv run python -m src.contracts.export_schemas --check  # anti-drift de los JSON Schema
\end{lstlisting}

\section{Construcción del corpus}
\label{sec:app-corpus}

El pipeline encadena descarga cruda, parseo, fragmentación y auditoría. La construcción con red y las validaciones posteriores son:

\begin{lstlisting}
uv run python scripts/build_corpus.py             # descarga + procesa + chunk (requiere red)
uv run python scripts/validate_mvp_corpus.py --strict
uv run python scripts/audit_corpus.py --strict    # debe terminar "Sin flags"
\end{lstlisting}

El reprocesado local a partir del \emph{raw} ya descargado (sin red) se realiza con \texttt{process\_mvp\_corpus.py}, útil cada vez que mejora el \emph{parser} o el fragmentador.

\section{Indexación densa}
\label{sec:app-indexacion}

Es una carga pesada; conviene lanzarla en el servidor y, si se dispone, acelerar la codificación en GPU:

\begin{lstlisting}
uv run python scripts/generate_dense_index.py --model e5-large-instruct
uv run python scripts/validate_dense_index.py <ruta-al-bundle>
\end{lstlisting}

\section{Evaluación}
\label{sec:app-evaluacion}

\begin{lstlisting}
# Recuperacion: comparacion de estrategias y seleccion de modelo de embeddings
uv run python scripts/benchmark_retrieval_strategies.py --bundle <bundle> --split test
uv run python scripts/benchmark_dense_models.py
# Generacion: retrieval + generacion + juez (opcional)
uv run python scripts/run_generation_eval.py
# Auditoria del banco de pruebas: grounding contra el corpus real
uv run python scripts/audit_eval_dataset.py       # debe terminar "Sin flags"
\end{lstlisting}

Cada corrida deja un \textbf{informe versionado e inmutable} en \texttt{data/processed/reports/}, con su configuración y sus huellas de reproducibilidad; de esos informes procede cada tabla del capítulo de experimentos.

\section{Ejecución real con modelos (opcional)}
\label{sec:app-ollama}

Las llamadas reales a \texttt{Ollama} se mantienen fuera de la suite automática. Para ejercitarlas de forma explícita:

\begin{lstlisting}
RUN_OLLAMA_INTEGRATION=1 uv run --locked pytest tests/test_integration_ollama.py -q -s
\end{lstlisting}

\chapter{Estructura del repositorio y licencias}
\label{app:repo}

\section{Enlace al repositorio}
\label{sec:app-enlace}

El código fuente completo de este Trabajo Fin de Grado está disponible en: \textcolor{red}{AÑADIR URL CUANDO SE PONGA PÚBLICO EL REPOSITORIO}.

\section{Estructura de carpetas}
\label{sec:app-estructura}

El repositorio se organiza en torno a una biblioteca (\texttt{src/}), sus \emph{scripts} de orquestación y los artefactos de datos y documentación:

\begin{description}
  \item[\texttt{src/}] La biblioteca del sistema, importable como \texttt{src.<paquete>}. Sus paquetes se corresponden uno a uno con las etapas del \emph{pipeline} y se describen en detalle en el capítulo de diseño (§\ref{subsec:mapa-componentes}): \texttt{boe}, \texttt{preprocessing}, \texttt{contracts}, \texttt{embeddings}, \texttt{indexing}, \texttt{retrieval}, \texttt{generation}, \texttt{evaluation}, \texttt{quality}, \texttt{config} y \texttt{core}.
  \item[\texttt{scripts/}] Los \emph{scripts} de orquestación de cada etapa (descarga, parseo, fragmentación, auditoría, indexación y evaluación).
  \item[\texttt{tests/}] La batería de pruebas, que se ejecuta por completo sin red mediante dobles de prueba.
  \item[\texttt{prompts/}] Las plantillas de \emph{prompt} versionadas (sistema, generación fundamentada y juez).
  \item[\texttt{schemas/}] Los \emph{JSON Schema} generados de forma determinista a partir de los contratos de datos.
  \item[\texttt{data/}] Los artefactos: \emph{raw} inmutable, procesados, índices e informes. En su mayoría no se versionan (sí los manifiestos, esquemas, informes de verificación y el banco de evaluación).
  \item[\texttt{thesis/}] Esta memoria, en \LaTeX.
  \item[\texttt{pyproject.toml}] La declaración de dependencias del proyecto, gestionadas por \texttt{uv} con fichero de bloqueo para un entorno reconstruible.
\end{description}

Esta correspondencia directa entre la estructura del código y las etapas de la cadena de datos es la que permite localizar cualquier pieza y reproducir cualquier resultado a partir de sus entradas.

\section{Licencias}
\label{sec:app-licencias}

El trabajo se apoya en materiales y componentes con licencias diferenciadas, todas compatibles con un uso abierto y local:

\begin{itemize}
  \item \textbf{Textos normativos.} Proceden de los datos abiertos del BOE y son reutilizables; tienen carácter informativo y no sustituyen a la publicación oficial \cite{boeLegislacionActualizada}.
  \item \textbf{Modelos.} El modelo de \emph{embeddings}, el generador y el juez son de pesos abiertos y se ejecutan en local, cada uno bajo la licencia que establece su autor.
  \item \textbf{Código.} Se publica como software de código abierto; la licencia concreta se incorpora al hacerse público el repositorio.
  \item \textbf{Memoria.} Se distribuye bajo licencia Creative Commons Reconocimiento-CompartirIgual (CC BY-SA 4.0).
\end{itemize}
